\section{Discussion}
% ============================================================

\subsection{Limitations}
\textit{[Placeholder: (1)~Assumes NL input is internally consistent; no contradiction
detection or resolution.
(2)~Multivariate distribution support is partial---see \S\ref{subsec:multivariate};
copula or conditional-chain design not yet finalised.
(3)~Deterministic inverse synthesis covers only single-dependency columns;
many-to-many joins and compound multi-column aggregations fall back to LLM.
(4)~Context enrichment relies on web retrieval quality; domain jargon absent from
public knowledge bases may still be under-specified.]}

\subsection{Failure Modes}
\textit{[Placeholder: (a)~Overly ambiguous NL where enrichment cannot recover missing
distribution parameters (e.g., no public definition of a proprietary scoring model).
(b)~DAG cycles introduced by contradictory logical constraints that the patch agent
cannot resolve within the retry budget.
(c)~Distribution parameter hallucination for rare distribution families
(e.g., generalised extreme value distributions).
(d)~Smoke-test failures that cannot be repaired by the code-generation agent
within 5 retries---currently causes a pipeline abort.]}

\subsection{Future Directions}
\textit{[Placeholder: (1)~Multivariate distributions: copula-based joint sampling to model
cross-column correlations beyond univariate priors; Gaussian copulas as a first step.
(2)~Contradiction detection and resolution in NL input before fact extraction.
(3)~Extending deterministic code generation to multi-dependency chains
(iterative inverse composition for chained aggregations).
(4)~Schema evolution: incrementally updating the generated schema and program as the
NL description changes, without full pipeline re-execution.
(5)~Support for non-relational output (document stores, graph databases).]}

